%
% process has one thread with two stacks
%
\chapter{调度}
\label{CH:SCHED}

任何操作系统运行的进程数都可能超过计算机的 CPU 数量，因此需要一种机制在进程之间
\indextext{多路复用（multiplexing）}
CPU。理想情况下，这种共享应对用户进程透明。一种常见的方法是为每个进程营造拥有独立虚拟 CPU 的假象，方法是将进程
多路复用到硬件 CPU 上。本章阐述 xv6 如何实现这种多路复用。

在继续阅读本章之前，请先阅读 {\tt kernel/proc.h}、
{\tt kernel/swtch.S}，
以及 {\tt kernel/proc.c} 中的
{\tt yield()}、
{\tt sched()}
和
{\tt schedule()}。

%%
\section{多路复用}
%%

Xv6 在两种情况下切换 CPU 以实现多路复用：
首先，当进程发起阻塞式系统调用（需要等待）时，例如 \lstinline{read} 或 \lstinline{wait}，xv6 会进行切换。
其次，xv6 定期强制切换，以处理长时间计算而不阻塞的进程。
前者称为自愿切换，
后者称为非自愿切换。

实现多路复用面临若干挑战。首先，如何从一个进程切换到另一个？基本思路是保存和恢复 CPU 寄存器，但这无法用 C 语言表达，因此变得棘手。其次，如何以对用户进程透明的方式强制切换？Xv6 采用标准技术，使用硬件定时器触发中断来驱动上下文切换。第三，所有 CPU 都在同一组进程之间切换，因此需要锁机制来避免错误，例如两个 CPU 同时决定运行同一个进程。第四，进程退出时必须释放其内存和其他资源，但进程自身无法完成所有这些工作。第五，多核机器的每个 CPU 必须记住它正在执行哪个进程，以便系统调用影响正确的进程内核状态。

%%
\section{上下文切换概述}
%%

术语``上下文切换（context switch）''指 CPU 停止执行一个内核线程（通常稍后恢复）并恢复执行另一个内核线程所涉及的步骤；这种切换是多路复用的核心。Xv6 不直接从一个进程的内核线程切换到另一个进程的内核线程；相反，内核线程通过切换到该 CPU 的``调度器线程（scheduler thread）''来放弃 CPU，然后调度器线程选择一个不同的进程内核线程来运行，并切换到该线程。

\begin{figure}[t]
\centering
\includegraphics[scale=0.5]{fig/switch.pdf}
\caption{从一个用户进程切换到另一个用户进程。在此示例中，xv6 运行在一个 CPU 上（因此有一个调度器线程）。}
\label{fig:switch}
\end{figure}

在更宏观的视角下，从一个用户进程切换到另一个用户进程所涉及的步骤如图~\ref{fig:switch} 所示：从旧进程的用户空间陷入（系统调用或中断）到其内核线程，切换到当前 CPU 的调度器线程，切换到新进程的内核线程，然后陷入返回到用户级进程。


%%
\section{代码：上下文切换}
%%

{\tt kernel/swtch.S} 中的函数 {\tt swtch()} 包含线程上下文切换的核心：它保存被切出线程的 CPU 寄存器，并恢复被切入线程先前保存的寄存器。这样足够的基本原因是线程状态由内存中的数据（如其栈）加上其 CPU 寄存器组成；线程内存无需保存和恢复，因为不同线程将数据保存在 RAM 的不同区域；但 CPU 只有一组寄存器，因此必须在线程之间切换（保存和恢复）它们。

% why must it be in assembler? why not C?
% why swtch(), not switch()?

每个线程的 {\tt struct proc} 包含一个 {\tt struct context}，用于在线程不运行时保存已保存的寄存器。CPU 调度器线程的 {\tt struct context} 位于该 CPU 的 {\tt struct cpu} 中。当线程 X 希望切换到线程 Y 时，线程 X 调用 {\tt swtch(\&X's context, \&Y's context)}。{\tt swtch()} 将当前 CPU 寄存器保存在 X 的上下文中，然后将 Y 的上下文内容加载到 CPU 寄存器中，然后返回。

以下是 {\tt swtch} 的简化副本：

\begin{verbatim}
swtch:
        sd ra, 0(a0)
        sd sp, 8(a0)
        sd s0, 16(a0)
        ...
        sd s11, 104(a0)

        ld ra, 0(a1)
        ld sp, 8(a1)
        ld s0, 16(a1)
        ...
        ld s11, 104(a1)

        ret
\end{verbatim}

{\tt a0} 保存第一个函数参数，{\tt a1} 保存第二个参数；在本例中，是两个 {\tt struct context} 指针。{\tt 16(a0)} 指从 {\tt a0} 指向的内存偏移 16 字节；参考 {\tt kernel/proc.h} 中 {\tt struct context} 的定义 \lineref{kernel/proc.h:/^struct.context/}，这是称为 {\tt s0} 的结构体字段。

{\tt swtch} 的 {\tt ret} 返回到哪里？它返回到 {\tt ra} 寄存器指向的指令。在线程 X 调用 {\tt swtch()} 切换到 Y 的示例中，当 {\tt ret} 执行时，{\tt ra} 刚从 Y 的 {\tt struct context} 加载。而 Y 的 {\tt struct context} 中的 {\tt ra} 最初是由 Y 调用 {\tt swtch} 时保存的，当时 Y 放弃了 CPU。因此 {\tt ret} 返回到 Y 调用 {\tt swtch()} 之后的指令；也就是说，X 调用 {\tt swtch()} 返回，就好像从 Y 最初调用 {\tt swtch()} 返回一样。而且 {\tt sp} 将是 Y 的栈指针，因为 {\tt swtch} 从 Y 的 {\tt struct context} 加载了 {\tt sp}；因此在返回时，Y 将在自己的栈上执行。{\tt swtch()} 无需直接保存或恢复程序计数器；保存和恢复 {\tt ra} 就足够了。

\lstinline{swtch}
\lineref{kernel/swtch.S:/swtch/}
保存被调用者保存的寄存器（{\tt ra,sp,s0..s11}），但不保存调用者保存的寄存器。
RISC-V 调用约定要求，如果代码需要在函数调用之间保留调用者保存寄存器中的值，
编译器必须生成指令，在函数调用之前将寄存器保存到栈中，
并在函数返回时从栈中恢复。
因此 \lstinline{swtch} 可以依赖调用它的函数已经保存了调用者保存的寄存器（要么是这样，
要么是调用函数不需要寄存器中的值）。

%%
\section{代码：调度}
%%

% why a separate thread

上一节介绍了
\indexcode{swtch}
的内部实现；
现在让我们将
\lstinline{swtch}
视为给定，并研究从一个进程的内核线程
通过调度器切换到另一个进程。
调度器以每个 CPU 一个特殊线程的形式存在，每个线程运行
\lstinline{scheduler}
函数。
该函数负责选择接下来运行哪个进程。
每个 CPU 都有自己的调度器线程，因为可能有多个 CPU 在任何给定时间寻找要运行的东西。
进程切换总是通过调度器线程进行，而不是直接从一个进程切换到另一个进程，
以避免在某些情况下调度器没有栈可以执行（例如，如果旧进程已退出，或当前没有其他进程想要运行）。

想要放弃 CPU 的进程必须
获取自己的进程锁
\indexcode{p->lock}，
释放它持有的任何其他锁，
更新自己的状态
(\lstinline{p->state})，
然后调用
\indexcode{sched}。
你可以在
\lstinline{yield}
\lineref{kernel/proc.c:/^yield/}、
\texttt{sleep}
和
\texttt{kexit}
中看到这一系列操作。
\indexcode{sched}
调用
\indexcode{swtch}
将当前上下文保存在
\lstinline{p->context}
中，并切换到
\indexcode{cpu->context}
中的调度器上下文。
\lstinline{swtch}
在调度器的栈上返回，
就好像
\indexcode{scheduler}
的
\lstinline{swtch}
已经返回一样
\lineref{kernel/proc.c:/swtch.*contex.*contex/}。

\lstinline{scheduler}
\lineref{kernel/proc.c:/^scheduler/}
运行一个循环：
找到一个要运行的进程，{\tt swtch()} 到它，最终它会 {\tt swtch()} 回调度器，然后调度器继续其循环。
调度器遍历进程表，
寻找可运行的进程，即具有
\lstinline{p->state}
\lstinline{==}
\lstinline{RUNNABLE}
的进程。
一旦找到进程，它就设置每个 CPU 的当前进程
变量
\lstinline{c->proc}，
将进程标记为
\lstinline{RUNNING}，
然后调用
\indexcode{swtch}
开始运行它
\linerefs{kernel/proc.c:/Switch.to/,/swtch/}。
在过去的某个时刻，目标进程必须调用了 {\tt swtch()}；
调度器对 {\tt swtch()} 的调用有效地从该早期调用返回。
图~\ref{fig:inout} 说明了这种模式。

\begin{figure}[t]
\input{fig/switch.tex}
\caption{{\tt swtch()} 总是以调度器线程作为源或目标，并且相关的 {\tt p->lock} 总是保持持有。}
\label{fig:inout}
\end{figure}

xv6 在调用
\lstinline{swtch}
时保持
\indexcode{p->lock}
：
\lstinline{swtch}
的调用者获取锁，
但在 {\tt swtch} 返回后在目标中释放它。
这种安排是不寻常的：更常见的是获取锁的线程也释放它。
Xv6 的上下文切换打破了这一约定，因为
\indexcode{p->state}
和
\lstinline{p->context}
必须一起原子地更新。
例如，如果在调用
\indexcode{swtch}
之前释放
\lstinline{p->lock}，
不同的 CPU $c$ 可能会决定运行该进程，因为
其状态是 \lstinline{RUNNABLE}。
CPU $c$ 将调用 \lstinline{swtch}，
它将从 \lstinline{p->context} 恢复，
而原始 CPU 仍在向
\lstinline{p->context}
保存。
结果将是该进程将在 CPU $c$ 上用部分保存的寄存器恢复，
而且两个 CPU 将使用相同的栈，这会导致混乱。
一旦 \lstinline{yield} 开始修改运行进程的状态
使其成为
\lstinline{RUNNABLE}，
\lstinline{p->lock} 必须保持持有，直到
进程已保存所有寄存器且
调度器在其栈上运行。
最早正确的释放点是在
\lstinline{scheduler}
（在其自己的栈上运行）
清除
\lstinline{c->proc}
之后。
同样，一旦
\lstinline{scheduler}
开始将 \lstinline{RUNNABLE} 进程转换为
\lstinline{RUNNING}，
在该进程的内核线程完全运行之前（在
\lstinline{swtch}
之后，例如在
\lstinline{yield}
中）不能释放锁。

有一种情况，调度器对
\indexcode{swtch}
的调用不会最终进入
\indexcode{sched}。
\lstinline{allocproc} 将新进程的上下文 {\tt ra}
寄存器设置为
\indexcode{forkret}
\lineref{kernel/proc.c:/^forkret/}，
因此其第一次 \lstinline{swtch} ``返回''到该函数的开头。
\lstinline{forkret}
存在是为了释放
\indexcode{p->lock}
并设置一些控制寄存器和陷入帧字段，
这些是返回用户空间所需的。
最后，{\tt forkret} 模拟从系统调用返回到用户空间的正常返回路径。


%%
\section{代码：mycpu 和 myproc}
%%

 xv6 经常需要获取指向当前进程的 \lstinline{proc} 结构体的指针。
在单核处理器上，可以使用全局变量指向当前 \lstinline{proc}。
这在多核机器上行不通，因为每个 CPU 执行不同的进程。
解决这个问题的方法是利用每个 CPU 都有自己的寄存器组这一事实。

当给定 CPU 在内核中执行时，xv6 确保 CPU 的 \lstinline{tp}
寄存器始终保存 CPU 的 hartid。
RISC-V 对其 CPU 进行编号，给每个一个唯一的 \indextext{hartid}。
\indexcode{mycpu}
\lineref{kernel/proc.c:/^mycpu/}
使用 \lstinline{tp} 索引一个
\lstinline{cpu} 结构体数组，并返回当前 CPU 的那个结构体。
\indexcode{struct cpu}
\lineref{kernel/proc.h:/^struct.cpu/}
保存一个指向正在该 CPU 上运行的
进程的 \lstinline{proc} 结构体的指针（如果有），
CPU 调度器线程的已保存寄存器，
以及管理中断禁用所需的嵌套自旋锁计数。

确保 CPU 的 \lstinline{tp} 保存 CPU 的 hartid 有点复杂，因为用户代码可以自由修改 \lstinline{tp}。
\lstinline{start} 在 CPU 启动序列的早期设置 \lstinline{tp}
寄存器，同时仍在机器模式中
\lineref{kernel/start.c:/w_tp/}。
在准备返回到用户空间时，
\lstinline{prepare_return} 将 \lstinline{tp} 保存在跳板页中，以防用户代码修改它。
最后，\lstinline{uservec} 在从用户空间进入内核时恢复保存的 \lstinline{tp}
\lineref{kernel/trampoline.S:/make tp hold/}。
编译器保证在内核代码中从不修改 \lstinline{tp}。
如果 xv6 可以在需要时询问 RISC-V 硬件获取当前 hartid 会更方便，
但 RISC-V 只允许在机器模式中这样做，而在监管者模式中不行。

\lstinline{cpuid}
和
\lstinline{mycpu}
的返回值是脆弱的：
如果定时器中断并导致线程让出，然后稍后在不同的 CPU 上恢复执行，
先前返回的值将不再正确。
为了避免这个问题，xv6 要求代码在调用 {\tt cpuid()} 或 {\tt mycpu()} 之前禁用中断，
并且只在使用完返回值后才启用中断。

函数
\indexcode{myproc}
\lineref{kernel/proc.c:/^myproc/}
返回当前 CPU 上运行的进程的
\lstinline{struct proc}
指针。
\lstinline{myproc}
禁用中断，调用
\lstinline{mycpu}，
从
\lstinline{struct cpu}
中获取当前进程指针
(\lstinline{c->proc})，
然后启用中断。
即使启用了中断，
\lstinline{myproc}
的返回值也是安全可用的：
如果定时器中断将调用进程移动到不同的 CPU，其
\lstinline{struct proc}
指针将保持不变。


%%
\section{现实世界}
%%

xv6 调度器实现了一个简单的调度策略，轮流执行每个进程。
这种策略叫做
\indextext{轮询（round robin）}。
真实操作系统实现了更复杂的策略，例如，允许进程具有优先级。
其思想是可运行的高优先级进程将被调度器优先于可运行的低优先级进程。
这些策略可能变得复杂，因为通常有相互竞争的目标：
例如，操作系统可能还想保证公平性和高吞吐量。

%%
\section{练习}
%%

\begin{enumerate}

\item 修改 xv6，使得从一个进程的内核线程切换到另一个进程时只使用一次上下文切换，
  而不是通过调度器线程切换。让出的线程需要自己选择下一个线程并调用 \texttt{swtch}。
  挑战将是防止多个 CPU 意外执行同一线程；
  正确处理锁定；以及避免死锁。

\end{enumerate}
